\begin{solapartat}
L'onada a l'estadi seria una ona \underline{mecànica} \punts{0,2} (necessita d'un medi per propagar-se), \underline{transversal} \punts{0,2} (el sentit d'oscil·lació i de propagació són perpendiculars) i bidimensional (si considerem que la grada és un pla).

\medskip
La longitud d'ona és \underline{$\lambda = 50\ \mathrm{m}$} \punts{0,2} (distància entre dos punts que es mouen exactament igual) i la seva pulsació és: $\omega = \dfrac{2\pi}{T} = \dfrac{2\pi}{10} = 0,2\pi = 0,63\ \mathrm{rad/s}$ \punts{0,4}
\end{solapartat}

\begin{solapartat}
Si l'espectador descriu un moviment harmònic simple, l'equació del moviment es pot escriure de la forma: $y(t) = A\sin(\omega t + \varphi_0)$ \punts{0,2} on:
\begin{itemize}
  \item $A$ és l'amplitud del moviment
  \item $\omega$ és la freqüència angular i
  \item $\varphi_0$ és la fase inicial
\end{itemize}
Segons les dades de l'enunciat, determinem les constants:

\medskip
\punts{0,2} $A = \dfrac{1}{2} = 0,5\ \mathrm{m}$

\medskip
\punts{0,2} $y(t = 0) = -A = -0,5\ \mathrm{m}$

\medskip
\punts{0,2} $y(t = 0) = A\sin\varphi_0 = -A \Rightarrow \sin\varphi_0 = -1 \Rightarrow \varphi_0 = -\dfrac{\pi}{2}\ \mathrm{rad}$

\medskip
\punts{0,2} $y(t) = 0,5\ \mathrm{m}\cdot\sin\left(0,2\pi\,\dfrac{\mathrm{rad}}{\mathrm{s}}\cdot t - \dfrac{\pi}{2}\ \mathrm{rad}\right)$

\medskip
Si escriuen l'equació com un cosinus és correcte i només canvia la fase inicial que seria $\pm\pi$
\end{solapartat}
